
\documentclass{article}
\usepackage{colortbl}
\usepackage{makecell}
\usepackage{multirow}
\usepackage{supertabular}

\begin{document}

\newcounter{utterance}

\centering \large Interaction Transcript for game `wordle', experiment `medium\_frequency\_words\_no\_clue\_no\_critic', episode 8 with Qwen3.5{-}4B.
\vspace{24pt}

{ \footnotesize  \setcounter{utterance}{1}
\setlength{\tabcolsep}{0pt}
\begin{supertabular}{c@{$\;$}|p{.15\linewidth}@{}p{.15\linewidth}p{.15\linewidth}p{.15\linewidth}p{.15\linewidth}p{.15\linewidth}}
    \# & \multicolumn{2}{c}{Player} && \multicolumn{2}{c}{Game Master} \\
    \hline

    \theutterance \stepcounter{utterance}  
    & & & \multicolumn{4}{p{0.6\linewidth}}{
        \cellcolor[rgb]{0.9,0.9,0.9}{
            \makecell[{{p{\linewidth}}}]{
                \texttt{\tiny{[P1$\langle$GM]}}
                \texttt{You are a language wizard who likes to guess words by using the given rules.} \\
\\ 
\texttt{Welcome to Wordle! You have six attempts to guess the target word, a valid English word of five lowercase letters (a{-}z). Please use the tags "explanation:" and "guess:" to provide a concise explanation for each guess.} \\
\\ 
\texttt{For instance, if your guess is "apple", your response should be} \\
\texttt{explanation: this is a common five{-}letter English word, and I am starting my guess with this word.} \\
\texttt{guess: apple} \\
\\ 
\texttt{After each guess, your answer will be validated, and you will receive feedback indicating which letters are correct (green), which letters are correct but in the wrong position (yellow), and which letters are incorrect (red). This feedback can be useful in determining which letters to include or exclude in your next guess.} \\
\\ 
\texttt{For example, the feedback for "apple" might be:} \\
\texttt{guess\_feedback: a<yellow> p<yellow> p<green> l<yellow> e<red>} \\
\\ 
\texttt{The explanation should contain details about how the guess\_feedback is used to arrive at a new guess.} \\
\\ 
\texttt{Do not generate any other text.} \\
\\ 
\texttt{Let's begin with your first guess.} \\
            }
        }
    }
     \\ \\

    \theutterance \stepcounter{utterance}  
    & & \multicolumn{4}{p{0.6\linewidth}}{
        \cellcolor[rgb]{0.95,0.95,0.95}{
            \makecell[{{p{\linewidth}}}]{
                \texttt{\tiny{[P1$\leftrightarrow$P]}}
                \texttt{{[}DRAFT{]}} \\
\texttt{explanation: I will start with a very common five{-}letter word that contains a variety of frequent letters to maximize information gain.} \\
\texttt{guess: apple} \\
            }
        }
    }
    & \\ \\

    \theutterance \stepcounter{utterance}  
    & & \multicolumn{4}{p{0.6\linewidth}}{
        \cellcolor[rgb]{0.95,0.95,0.95}{
            \makecell[{{p{\linewidth}}}]{
                \texttt{\tiny{[P1$\leftrightarrow$P]}}
                \texttt{{[}SELF{-}FEEDBACK{]}} \\
\texttt{FORMAT: ok} \\
\texttt{RULES: ok} \\
\texttt{STRATEGY: ok} \\
\texttt{REASON: The response correctly uses the required tags and markers, follows the game's constraints, and selects a statistically strong opening word for Wordle.} \\
            }
        }
    }
    & \\ \\

    \theutterance \stepcounter{utterance}  
    & \multicolumn{4}{p{0.6\linewidth}}{
        \cellcolor[rgb]{0.561,0.737,0.561}{
            \makecell[{{p{\linewidth}}}]{
                \texttt{\tiny{[P1$\rangle$GM]}}
                \texttt{explanation: I will start with a very common five{-}letter word that contains a variety of frequent letters to maximize information gain.} \\
\texttt{guess: apple} \\
            }
        }
    }
    & & \\ \\

    \theutterance \stepcounter{utterance}  
    & & & \multicolumn{2}{p{0.3\linewidth}}{
        \cellcolor[rgb]{0.9,0.9,0.9}{
            \makecell[{{p{\linewidth}}}]{
                \texttt{\tiny{[GM$|$GM]}}
                \texttt{attempts = 1} \\
\texttt{target\_word = renew} \\
\texttt{guess = apple} \\
\texttt{guess\_feedback = a<red> p<red> p<red> l<red> e<yellow>} \\
            }
        }
    }
    & & \\ \\

    \theutterance \stepcounter{utterance}  
    & & & \multicolumn{4}{p{0.6\linewidth}}{
        \cellcolor[rgb]{0.9,0.9,0.9}{
            \makecell[{{p{\linewidth}}}]{
                \texttt{\tiny{[P1$\langle$GM]}}
                \texttt{guess\_feedback: a<red> p<red> p<red> l<red> e<yellow>} \\
\\ 
\texttt{Provide your response only in this format.} \\
\texttt{explanation: details} \\
\texttt{guess: word} \\
            }
        }
    }
     \\ \\

    \theutterance \stepcounter{utterance}  
    & & \multicolumn{4}{p{0.6\linewidth}}{
        \cellcolor[rgb]{0.95,0.95,0.95}{
            \makecell[{{p{\linewidth}}}]{
                \texttt{\tiny{[P1$\leftrightarrow$P]}}
                \texttt{{[}DRAFT{]}} \\
\texttt{explanation: The feedback indicates that 'e' is in the word but in the last position, while 'a', 'p', and 'l' are not present. I need to find a word ending in 'e' that does not contain 'a', 'p', or 'l'. I will try a common word ending in 'e' with different vowels and consonants.} \\
\texttt{guess: crane} \\
            }
        }
    }
    & \\ \\

    \theutterance \stepcounter{utterance}  
    & & \multicolumn{4}{p{0.6\linewidth}}{
        \cellcolor[rgb]{0.95,0.95,0.95}{
            \makecell[{{p{\linewidth}}}]{
                \texttt{\tiny{[P1$\leftrightarrow$P]}}
                \texttt{{[}SELF{-}FEEDBACK{]}} \\
\texttt{FORMAT: ok} \\
\texttt{RULES: ok} \\
\texttt{STRATEGY: problem} \\
\texttt{REASON: The guess "crane" is invalid because the feedback confirmed the letter 'e' is in the word, so a second guess cannot contain 'e' in any position.} \\
            }
        }
    }
    & \\ \\

    \theutterance \stepcounter{utterance}  
    & \multicolumn{4}{p{0.6\linewidth}}{
        \cellcolor[rgb]{0.561,0.737,0.561}{
            \makecell[{{p{\linewidth}}}]{
                \texttt{\tiny{[P1$\rangle$GM]}}
                \texttt{explanation: The feedback indicates that 'e' is in the word but in the last position, while 'a', 'p', and 'l' are not present. I need to find a word ending in 'e' that does not contain 'a', 'p', or 'l'. I will try a common word ending in 'e' with different vowels and consonants.} \\
\texttt{guess: crane} \\
            }
        }
    }
    & & \\ \\

    \theutterance \stepcounter{utterance}  
    & & & \multicolumn{2}{p{0.3\linewidth}}{
        \cellcolor[rgb]{0.9,0.9,0.9}{
            \makecell[{{p{\linewidth}}}]{
                \texttt{\tiny{[GM$|$GM]}}
                \texttt{attempts = 2} \\
\texttt{target\_word = renew} \\
\texttt{guess = crane} \\
\texttt{guess\_feedback = c<red> r<yellow> a<red> n<yellow> e<yellow>} \\
            }
        }
    }
    & & \\ \\

    \theutterance \stepcounter{utterance}  
    & & & \multicolumn{4}{p{0.6\linewidth}}{
        \cellcolor[rgb]{0.9,0.9,0.9}{
            \makecell[{{p{\linewidth}}}]{
                \texttt{\tiny{[P1$\langle$GM]}}
                \texttt{guess\_feedback: c<red> r<yellow> a<red> n<yellow> e<yellow>} \\
\\ 
\texttt{Provide your response only in this format.} \\
\texttt{explanation: details} \\
\texttt{guess: word} \\
            }
        }
    }
     \\ \\

    \theutterance \stepcounter{utterance}  
    & & \multicolumn{4}{p{0.6\linewidth}}{
        \cellcolor[rgb]{0.95,0.95,0.95}{
            \makecell[{{p{\linewidth}}}]{
                \texttt{\tiny{[P1$\leftrightarrow$P]}}
                \texttt{{[}DRAFT{]}} \\
\texttt{explanation: The feedback shows that 'r' and 'n' are in the word but not in the second and fourth positions respectively, and 'e' is in the word but not in the last position (contradicting my previous guess "crane" where I assumed 'e' was last; actually, the previous feedback for "crane" showed e<yellow>, meaning 'e' is present but not in the 5th spot. Wait, let me re{-}evaluate the previous turn. In "crane", 'e' was at the end and got yellow, so 'e' is in the word but not at the end. My previous deduction was slightly off. Now, with "crane" giving r<yellow> (pos 2), n<yellow> (pos 4), e<yellow> (pos 5), I know the word contains r, n, e. The positions for r and n are wrong. I need a word with r, n, e where they are not in those spots. Let's try a word like "range" but that ends in e. Or "renne"? No. Let's try "range" again? No, 'a' was red in "crane". So no 'a'. Let's try "sneer"? No 'r' at pos 2. How about "trench"? Too long. Let's try "renne" is not a word. Maybe "rune"? No. Let's try "sneak"? No 'e' at end. Let's try "sneer" again? 'e' at pos 2 and 4. Let's try "sneer". Wait, 'e' is yellow in "crane" at pos 5. So 'e' is not at pos 5. In "sneer", 'e' is at 2 and 4. Let's try "sneer". But I need to confirm 'r' and 'n'. Let's try "sneer". Actually, let's try "renne" is not valid. How about "rune" is 4 letters. Let's try "sneer". Wait, I need a 5 letter word. "sneer" has e at 2 and 4. If the target is "sneer", then in "crane", 'r' is at 2 (yellow, so not 2), 'n' is at 4 (yellow, so not 4), 'e' is at 5 (yellow, so not 5). This fits "sneer" if 'r' is at 3 or 4? No, 'r' is at 3 in sneer. 'n' is at 2. 'e' is at 1, 3, 4, 5? No. "s} \\
            }
        }
    }
    & \\ \\

    \theutterance \stepcounter{utterance}  
    & & \multicolumn{4}{p{0.6\linewidth}}{
        \cellcolor[rgb]{0.95,0.95,0.95}{
            \makecell[{{p{\linewidth}}}]{
                \texttt{\tiny{[P1$\leftrightarrow$P]}}
                \texttt{{[}SELF{-}FEEDBACK{]}} \\
\texttt{FORMAT: problem} \\
\texttt{RULES: problem} \\
\texttt{STRATEGY: problem} \\
\texttt{REASON: The response fails to follow the required output format, contains excessive internal monologue instead of a concise explanation, and proposes an invalid 4{-}letter word ("rune") as a guess.} \\
            }
        }
    }
    & \\ \\

    \theutterance \stepcounter{utterance}  
    & \multicolumn{4}{p{0.6\linewidth}}{
        \cellcolor[rgb]{0.561,0.737,0.561}{
            \makecell[{{p{\linewidth}}}]{
                \texttt{\tiny{[P1$\rangle$GM]}}
                \texttt{explanation: The feedback from "crane" indicates that 'r', 'n', and 'e' are in the word, but 'r' is not in the second position, 'n' is not in the fourth position, and 'e' is not in the fifth position. I need a five{-}letter word containing r, n, and e with these letters in different positions. "sneer" fits the criteria: it contains s, n, e, e, r. In "sneer", 'n' is in the second position (which contradicts the yellow 'n' in "crane" at position 4? Wait, in "crane", 'n' was at position 4 and got yellow, meaning 'n' is in the word but not at position 4. In "sneer", 'n' is at position 2. This is consistent. 'r' in "crane" was at position 2 and got yellow, so 'r' is not at position 2. In "sneer", 'r' is at position 5. This is consistent. 'e' in "crane" was at position 5 and got yellow, so 'e' is not at position 5. In "sneer", 'e' is at positions 2 and 4. This is consistent. I will guess "sneer" to test the positions of these letters.} \\
\texttt{guess: sneer} \\
            }
        }
    }
    & & \\ \\

    \theutterance \stepcounter{utterance}  
    & & & \multicolumn{2}{p{0.3\linewidth}}{
        \cellcolor[rgb]{0.9,0.9,0.9}{
            \makecell[{{p{\linewidth}}}]{
                \texttt{\tiny{[GM$|$GM]}}
                \texttt{attempts = 3} \\
\texttt{target\_word = renew} \\
\texttt{guess = sneer} \\
\texttt{guess\_feedback = s<red> n<yellow> e<yellow> e<green> r<yellow>} \\
            }
        }
    }
    & & \\ \\

    \theutterance \stepcounter{utterance}  
    & & & \multicolumn{4}{p{0.6\linewidth}}{
        \cellcolor[rgb]{0.9,0.9,0.9}{
            \makecell[{{p{\linewidth}}}]{
                \texttt{\tiny{[P1$\langle$GM]}}
                \texttt{guess\_feedback: s<red> n<yellow> e<yellow> e<green> r<yellow>} \\
\\ 
\texttt{Provide your response only in this format.} \\
\texttt{explanation: details} \\
\texttt{guess: word} \\
            }
        }
    }
     \\ \\

    \theutterance \stepcounter{utterance}  
    & & \multicolumn{4}{p{0.6\linewidth}}{
        \cellcolor[rgb]{0.95,0.95,0.95}{
            \makecell[{{p{\linewidth}}}]{
                \texttt{\tiny{[P1$\leftrightarrow$P]}}
                \texttt{{[}DRAFT{]}} \\
\texttt{explanation: The feedback for "sneer" shows that 'e' is green (correct position), 'n' and 'r' are yellow (correct letters, wrong positions), and 's' is red (not in the word). This means the word contains two 'e's, one 'n', and one 'r'. The 'e' must be in the third position. The 'n' is not in the second position, and the 'r' is not in the fifth position. I need to find a word with the pattern \_ \_ e e \_ where the second letter is not 'n' and the fifth letter is not 'r'. "beers" fits this pattern perfectly: 'b' is new, 'e' is green, 'e' is green, 'r' is yellow (not at the end), 's' is new. Wait, "beers" has 's' at the end. Let's re{-}evaluate. The word has two 'e's. One is at pos 3. The other 'e' is likely at pos 4 based on "sneer" giving two yellows for 'e'? No, "sneer" gave one green 'e' at pos 3 and one yellow 'e' at pos 4. So the word has 'e' at 3 and 'e' at 4. The 'n' is not at 2. The 'r' is not at 5. The word is \_ \_ e e \_. "beers" has 'r' at 4 and 's' at 5. If the target is "beers", then 'r' would be green at 4, but in "sneer" 'r' was at 5 and got yellow. This implies 'r' is in the word but not at 5. In "beers", 'r' is at 4. This is consistent. 's' is at 5 in "beers". In "sneer", 's' was at 1 and got red. So 's' is not in the word. Therefore "beers" is incorrect because it contains 's'. I need a word like \_ \_ e e \_ without 's'. "feels" has 's'. "peels" has 's'. "keels" has 's'. "neels" has 's'. "bees" is 4 letters. "crees" is not a word. "brees" is not a word. "treys"? No. Let's reconsider "sneer". 'n' was yellow at pos 2. So 'n' is in the word but not at pos 2. 'r' was yellow at pos 5. So 'r' is in the word but not at pos 5. 'e' was green at pos 3} \\
            }
        }
    }
    & \\ \\

    \theutterance \stepcounter{utterance}  
    & & \multicolumn{4}{p{0.6\linewidth}}{
        \cellcolor[rgb]{0.95,0.95,0.95}{
            \makecell[{{p{\linewidth}}}]{
                \texttt{\tiny{[P1$\leftrightarrow$P]}}
                \texttt{{[}SELF{-}FEEDBACK{]}} \\
\texttt{FORMAT: problem} \\
\texttt{RULES: problem} \\
\texttt{STRATEGY: problem} \\
\texttt{REASON: The response fails to follow the required output format, contains a hallucinated word ("beers"), and proposes an invalid guess ("treys") that violates the five{-}letter constraint.} \\
            }
        }
    }
    & \\ \\

    \theutterance \stepcounter{utterance}  
    & \multicolumn{4}{p{0.6\linewidth}}{
        \cellcolor[rgb]{0.561,0.737,0.561}{
            \makecell[{{p{\linewidth}}}]{
                \texttt{\tiny{[P1$\rangle$GM]}}
                \texttt{explanation: The feedback for "sneer" indicates that 'e' is green at position 3, 'n' is yellow (present but not at position 2), 'e' is yellow at position 4 (meaning the second 'e' is actually at position 4), and 'r' is yellow (present but not at position 5). 's' is red, so it is not in the word. The pattern is \_ \_ e e \_. I need a word with 'e' at positions 3 and 4, containing 'n' and 'r' elsewhere, but not 's'. "feels" is invalid because it has 's'. "peels" is invalid because it has 's'. "keels" is invalid because it has 's'. "neels" is invalid because it has 's'. "bees" is only four letters. "crees" is not a word. "brees" is not a word. "treys" is invalid because it has 's' and 'y'. "reels" fits the pattern: 'r' is at position 1 (consistent with being yellow in "sneer" at position 5), 'e' is at 2 (consistent with being yellow in "sneer" at position 4? No, wait. In "sneer", 'e' at pos 4 was yellow, meaning 'e' is in the word but not at pos 4. But in "reels", 'e' is at pos 2 and 4. If the target is "reels", then 'e' at pos 4 should be green. But in "sneer", 'e' at pos 4 was yellow. This implies the second 'e' is NOT at position 4. Let's re{-}read the feedback for "sneer": s<red> n<yellow> e<yellow> e<green> r<yellow>. Wait, the feedback says e<green> at pos 4? No, the feedback was s<red> n<yellow> e<yellow> e<green> r<yellow>. Let's re{-}parse the previous turn. "sneer": s<red> n<yellow> e<yellow> e<green> r<yellow>. This means: s is not in word. n is in word, not at pos 2. e is in word, not at pos 3. e is in word, at pos 4. r is in word, not at pos 5. So the pattern is \_ \_ \_ e \_. And we know 'e' is in the word. Wait, the previous turn said "e<green>" at pos 4? Let's check the feedback string again: "s<red> n<yellow> e<yellow> e<green> r<yellow>". Yes} \\
            }
        }
    }
    & & \\ \\

    \theutterance \stepcounter{utterance}  
    & & & \multicolumn{2}{p{0.3\linewidth}}{
        \cellcolor[rgb]{0.9,0.9,0.9}{
            \makecell[{{p{\linewidth}}}]{
                \texttt{\tiny{[GM$|$GM]}}
                \texttt{Error: The response should contain the 'guess:' keyword exactly once.} \\
            }
        }
    }
    & & \\ \\

    \theutterance \stepcounter{utterance}  
    & & & \multicolumn{2}{p{0.3\linewidth}}{
        \cellcolor[rgb]{0.9,0.9,0.9}{
            \makecell[{{p{\linewidth}}}]{
                \texttt{\tiny{[GM$|$GM]}}
                \texttt{game\_result = ABORT} \\
            }
        }
    }
    & & \\ \\

\end{supertabular}
}

\end{document}
